% 面向动物机器人自动导航的地面站-边缘端协同控制架构研究


\chapter{绪论}
\section{研究背景与意义}
动物机器人Animal Robot),又称赛博动物Cyborg Animal)或生物混合机器人
Biohybrid Robot),是一类将微电子神经接口植入活体动物神经系统,通过外部
控制指令调控动物行为,从而实现特定导航或探测任务的新型智能体。
与传统机械机器人相比,动物机器人天然具备生物感知能力强、运动适应性高、
能量利用效率优越等显著优势,在灾害搜救、反恐排爆、地形探测等复杂非结构化
场景中展现出广阔的应用前景,因此自本世纪初以来持续受到机器人学、神经工程学
与智能控制领域的广泛关注\cite{talwar2002rat}。
 
在动物机器人的典型实现方案中,大鼠被认为是最具代表性的载体之一。
其神经系统结构相对清晰,电刺激效果稳定且可重复,已成为该领域最为深入研究的
对象\cite{zheng2013rat}。
浙江大学郑筱祥团队自2007年起长期深耕大鼠机器人研究,先后提出了基于视觉线索
引导的自动导航方案\cite{wang2015visual}、基于有限状态机的自主搜索系统
\cite{zheng2023fsm},以及近年来提出的无人机引导大鼠协同导航
Rat-UAV Navigation)范式\cite{zheng2024ratuav},逐步将大鼠机器人从简单
遥控推进至自主导航阶段。
 
然而,上述研究在系统架构层面存在一个共性局限：现有大鼠机器人控制系统普遍
沿用"个人计算机PC)集中处理\,+\,蓝牙点对点通信"的单体架构模式。
在这一架构下,感知、决策与控制指令下发均由地面端单台主机承担,大鼠背包端
仅执行简单的刺激输出；蓝牙通信的往返延迟通常在100至500毫秒之间
\cite{zhang2019bci},有效通信距离普遍不超过5至10米\cite{feng2007remote}。
上述约束在简单实验室场景中尚可接受,但在Rat-UAV协同导航这一更具现实意义的
任务模式下,无人机UAV)飞行范围远超蓝牙覆盖半径,链路质量随飞行距离
和障碍物遮挡产生剧烈波动,原有单体架构的实时性与可靠性均难以满足任务需求,
成为制约动物机器人走向实际部署的核心瓶颈。
 
与此同时,分布式边缘计算与微服务架构在机器人领域的快速发展为上述问题的
解决提供了新的技术路径。以云机器人Cloud Robotics)\cite{kehoe2015cloud}
和边缘加速机器人导航Edge Accelerated Robot Navigation,EARN)
\cite{li2025earn}为代表的研究表明,将计算密集型任务卸载至高算力节点,
同时在本地边缘端保留实时控制逻辑,可以在资源受限场景下有效降低端到端
控制延迟并提升系统可扩展性。然而,这些方案主要面向工业机器人或无人机等
确定性执行器,尚未针对动物机器人背包端带宽极窄、生物行为不确定性高等
特殊约束进行适应性设计,因此无法直接迁移至大鼠机器人控制场景。
 
基于上述背景,本文以浙江大学Rat-UAV协同导航范式为应用背景,针对动物机器人
自动导航任务中控制系统架构与通信协议两个核心维度的空白,开展面向地面站-边缘端
协同控制的分布式架构研究。本文的研究成果不仅能够为动物机器人实验提供更高性能
的控制平台,也将为分布式系统理论在生物混合机器人领域的工程化落地提供参考。


\section{国内外研究现状}
本节从三个相互独立且共同构成分布式协同控制系统完整图景的维度,系统梳理
国内外相关研究进展：首先介绍动物机器人控制任务的背景与典型系统,继而综述
分布式机器人控制架构的主要范式与演进脉络,最后分析实时通信协议的研究现状,
为后续研究问题的提出奠定文献基础。
\subsection{动物机器人控制任务背景}
动物机器人研究的起点通常追溯至2002年Talwar等人在《自然》杂志发表的开创性
工作\cite{talwar2002rat}：通过向大鼠体感皮层和内侧前脑束植入微电极,
利用无线遥控实现了活体大鼠在复杂地形中的可控导航,首次证明了外部电信号
干预动物行为的可行性,开创了"Robo-Rat"研究方向。
此后,以蟑螂\cite{ariyanto2024cockroach}、熊蜂\cite{PlaceholderBumbleBee}、
信鸽等为载体的多种动物机器人相继被报道,展现了该研究方向的广泛适用性。
 
在大鼠机器人自主导航方面,浙江大学团队的系列工作代表了国内最高水平。
Wang等人于2015年提出视觉线索引导框架,将摄像头采集的图像经无线传输至
PC端进行目标检测,再将结果转化为刺激序列闭环控制大鼠转向\cite{wang2015visual};
Zheng等人于2023年提出基于有限状态机的大鼠自主导航系统,集成IMU与红外传感器,
在开放场地实现无辅助摄像头的自主搜索\cite{zheng2023fsm},并在论文中明确
指出"目前缺乏构建动物机器人自主导航系统的总体框架"。
这一论断直接揭示了系统架构层面的研究空白,是本文选题的重要动机之一。
最新的Rat-UAV Navigation工作进一步引入无人机作为移动观测平台,通过POMDP
框架处理大鼠行为不确定性\cite{zheng2024ratuav},但其软件架构仍为PC集中式
模式,边缘端协同尚未实现。
 
综上,现有动物机器人研究在控制算法层面已取得显著进展,但在控制系统架构方面,
自2002年至今普遍停留于PC集中处理\,+\,蓝牙点对点的单体模式,尚无面向分布式
协同的系统架构设计。这一空白构成本文的核心研究问题。

\subsection{分布式机器人控制系统研究}
分布式机器人控制系统的架构范式经历了从集中式、云机器人到边缘协同、再到
微服务化的多代演进。Medeiros等人的经典综述将移动机器人控制架构归纳为
反应式、审议式和混合式三大类\cite{medeiros1998survey},奠定了控制架构
研究的基础分类框架。随着网络技术的发展,Kehoe等人提出云机器人范式,将感知、
规划等高算力任务卸载至云端,机器人本体仅保留执行功能\cite{kehoe2015cloud},
显著提升了计算能力上限,但公有云的网络往返延迟通常超过50\,ms)限制了
其在实时控制场景中的应用。
 
针对云端延迟过高的问题,边缘计算范式被引入机器人控制领域。
FogROS2是这一方向的代表性工作,Ichnowski等人基于ROS\,2实现了计算节点向
云端或雾节点的透明迁移,在SLAM等任务中将延迟降低约50\%\cite{ichnowski2023fogros2}。
Li等人提出的EARN框架在机器人本地与边缘端之间动态切换运动规划器,在
Jetson Nano与Orin NX上验证了边缘加速的有效性\cite{li2025earn},
但其侧重于规划算法切换,缺乏微服务架构与通信协议层面的系统性设计。
针对UAV等资源受限飞行平台,Satpure等人将MPC控制器基于
Kubernetes\,+\,Docker容器化部署于边缘服务器,使控制频率从机载20\,Hz
提升至100\,Hz\cite{satpure2022mpc},验证了容器化在机器人边缘控制中的可行性。
 
在微服务架构方面,Xia等人首次将微服务引入云机器人系统,利用Docker容器将
机器人功能分解为协作微服务并通过Kubernetes管理\cite{xia2018microservice}；
Schrick等人将微服务架构应用于安全关键移动机器人控制,证明微服务化可显著
提升系统并发性与容错能力\cite{schrick2024safety}。
然而,上述工作均在云端或资源充足的工业平台上实现,未涉及动物机器人背包端
资源极度受限通信带宽约200\,kbps、计算能力有限)的特殊场景,其方案无法
直接迁移。
 
综合来看,分布式机器人控制系统研究已在云机器人、边缘协同和微服务化三个
维度积累了丰富成果,但现有工作普遍面向工业机器人、无人机等确定性执行器,
其架构设计均未考虑生物机器人特有的不确定性和超低带宽约束。将上述分布式
架构理论适配至动物机器人场景,是本文需要重点解决的问题。

\subsection{实时通信协议在分布式控制中的研究}
通信协议是分布式机器人控制系统的神经网络,其延迟特性直接决定了系统的
控制实时性上限。现有机器人通信协议主要可分为三类：基于发布-订阅模型的
DDS/ROS\,2、基于请求-响应模型的gRPC/REST,以及面向IoT的MQTT等轻量级
协议\cite{profanter2019opcua}。
 
ROS\,2采用DDSData Distribution Service)作为中间件,其去中心化点对点
通信架构在理论上具有良好的实时性。然而,Maruyama等人的实测表明,ROS\,2
相比底层DDS可产生约50\%的延迟附加开销\cite{maruyama2016exploring}；
后续研究进一步指出,标准ROS\,2无法在嵌入式设备上保证1\,ms通信周期
\cite{PlaceholderROS2Realtime2023}。
这一特性使ROS\,2难以满足动物机器人小于11\,ms端到端延迟的严苛要求。
 
% gRPC凭借HTTP/2多路复用与Protocol\,Buffers高效序列化,在微服务通信场景中
% 展现出显著的延迟优势。对比测试表明,gRPC响应时间比REST\,API快约77\%
% \cite{PlaceholderGrpcRest2024}；Protocol\,Buffers序列化后消息体积约为
% JSON的$\frac{1}{2}$至$\frac{1}{3}$\cite{popic2016protobuf}。
% 然而,现有gRPC性能研究均在Web微服务或数据库读写场景中开展,缺乏在嵌入式
% 设备上、面向实时控制场景的系统性评估,也未涉及针对超窄带宽和远距离帧丢失
% 的协议级保护机制。
 
在动物机器人通信实践方面,现有系统普遍采用蓝牙作为地面站与背包端之间的
通信链路。Zhang等人的脑-脑接口大鼠控制实验实测显示,蓝牙控制指令延迟达
155至495\,ms\cite{zhang2019bci}；BLE在工业IoT场景中的评估同样表明,
单帧传输延迟约为10\,ms,大数据包传输延迟可达400\,ms至1\,s,且无法保证
最大传输延迟的确定性上界\cite{PlaceholderBLE2017}。
上述数据定量揭示了蓝牙方案在实时性和可靠性方面的根本局限,为本文引入基于
gRPC扩展协议的替代方案提供了充分动机。
 
综上所述,gRPC协议在延迟性能和序列化效率上具有显著优势,但其标准实现
存在协议头开销偏重、缺乏控制帧优先级保护和坏帧恢复机制等不足,尚未针对
动物机器人远距离导航的超窄带宽场景进行适应性扩展。这一研究空白是本文
提出mice-gRPC协议的核心动机。

\section{本文主要研究内容}

基于上述研究现状的梳理,本文面向动物机器人自动导航任务,以浙江大学
Rat-UAV协同导航范式为应用背景,针对地面站-边缘端协同控制架构和远距离
低延迟通信协议两个核心问题展开研究。本文的主要研究内容如下：
 
\begin{enumerate}[leftmargin=2em, label=(\arabic*)]
 
  \item \textbf{面向动物机器人的三层异构协同控制架构设计。}
  在Rat-UAV协同范式基础上,设计并实现"地面站高算力决策)$-$UAV边缘端
  实时感知与中继)$-$大鼠背包端刺激执行)"三层异构分布式架构,
  以微服务形式部署Control Service与Stream Service,通过Docker多架构容器化
  实现x86地面站与ARM64 Jetson Orin边缘节点的跨平台统一管理,
  为两个核心创新点提供运行基础。
 
  \item \textbf{面向弱网环境的低时延任务调度算法WLTS)。}
  针对UAV飞行过程中链路质量波动导致任务卸载延迟不可预测的问题,提出
  基于链路质量感知的静态-动态混合任务调度算法。算法将实时控制任务固定
  保留于UAV边缘端执行,计算密集型任务视觉处理、路径规划)依据链路
  健康度三级状态机动态决策卸载至地面站或本地降级执行,并引入容错降级
  机制保证弱网条件下控制不中断。
 
  \item \textbf{面向远距离导航的mice-gRPC通信协议。}
  针对大鼠背包端蓝牙通信带宽极窄约200\,kbps)、远距离传输帧丢失率高、
  标准gRPC协议头开销与控制帧抢占问题,在gRPC协议基础上提出
  mice-gRPCMinimal IoT Control Extended gRPC)扩展协议。
  协议创新点包括：针对控制指令设计的轻量化Protobuf帧格式
  目标帧大小小于20字节)、CRC16坏帧校验与ARQ优先级重传机制,
  以及控制指令与视频流的单通路时分复用调度算法。
 
\end{enumerate}
 
本文的两项核心创新点概括如下：
 
\begin{itemize}[leftmargin=2em]
 
  \item \textbf{创新点一：}
  提出面向弱网环境的低时延任务调度算法WLTS),首次将链路自适应任务调度
  机制引入动物机器人分布式控制系统,实现弱网条件下端到端控制延迟小于11\,ms,
  相较于现有蓝牙集中式方案100至500\,ms)提升90\%以上。
 
  \item \textbf{创新点二：}
  提出面向远距离导航的mice-gRPC通信协议,在标准gRPC基础上扩展轻量化帧
  设计、坏帧保护与单通路复用机制,支持导航距离突破原有蓝牙方案的3至5米
  限制、扩展至15米以上,为动物机器人远距离自主导航提供可靠通信保障。
 
\end{itemize}
\section{本文结构}

aaaaaaaaaaaaaaaaaaaa


\chapter{第2章 相关工作与技术综述}
实现一个面向动物机器人的地面站-边缘端分布式协同控制系统,需要在架构设计、任务调度与通信协议三个核心维度上提供系统性的理论支撑。三者相互依存：合理的分布式架构是任务调度得以运行的物理基础,高效的任务调度算法决定了架构的实时性能上限,而通信协议则是连接各层节点、保障指令可靠传输的底层神经。脱离任一维度的孤立研究均无法有效解决动物机器人远距离自主导航中的核心问题。
基于上述认识,本章围绕三类核心技术问题展开文献综述与理论分析。2.2节系统梳理分布式控制架构的设计方法,重点介绍从集中式到微服务化的范式演进、任务卸载理论与嵌入式微服务部署实践；2.3节从协议选型角度分析实时通信协议的性能特征,重点阐述gRPC协议原理与面向资源受限环境的优化方向；2.4节以问题驱动的视角剖析现有动物机器人信息传输系统的共性架构与量化局限,论证引入分布式协同架构的必要性。2.5节对本章内容作综合性小结,并指出现有研究存在的空白。

\section{分布式控制架构的设计方法}
本节从机器人控制架构的历史演进出发,依次介绍集中式到微服务化的范式转变、边缘协同架构的任务卸载理论,以及面向嵌入式设备的微服务部署实践,为本文三层异构协同控制架构的设计提供理论依据。

\subsection{机器人控制架构的范式演进}
机器人控制架构的研究伴随着计算机技术和网络技术的发展经历了显著的范式演进,从早期的集中式单体架构,逐步向分布式、云化与微服务化方向演进。理解这一演进脉络,是把握本文架构设计取舍的前提。
集中式架构与分布式架构的分野。 Medeiros 等人在经典综述中将移动机器人控制架构归纳为反应式Reactive)、审议式Deliberative)和混合式Hybrid)三大范式 [Medeiros et al., 1998]。早期机器人系统多采用集中式审议架构：单台主机承担感知、决策与控制指令下发的全部职责,结构简单、调试便捷,但随着任务复杂度的提升,集中式架构的单点瓶颈日益凸显——计算资源上限受限于单机算力,通信链路的故障即导致系统全局失效,且难以支持多机协同场景下的并发控制需求。
云机器人范式的兴起。 随着云计算基础设施的普及,Kehoe 等人系统提出了云机器人Cloud Robotics)范式：将感知数据处理、路径规划、模型推理等计算密集型任务卸载至云端服务器,机器人本体仅保留执行层功能,从而突破单机算力上限 [Kehoe et al., 2015]。云机器人范式显著提升了系统可扩展性,并为多机器人协同提供了统一的数据共享平台。然而,云端数据中心与机器人本体之间的网络往返延迟通常超过50 ms,在实时控制场景中难以满足严苛的时序要求,这一固有局限催生了边缘计算在机器人领域的应用需求。
边缘协同与雾计算范式。 以FogROS2为代表的研究将计算节点从公有云下沉至网络边缘,显著缩短了控制环路延迟。Ichnowski 等人基于 ROS 2 实现了计算节点向云端或雾节点的透明迁移,在 SLAM 和抓取规划等任务中延迟降低约50\%,运动规划加速达28至45倍 [Ichnowski et al., 2023]。与此同时,Macenski 等人正式发布 ROS 2 系统设计论文,系统阐述了其基于 DDS 的去中心化点对点通信、QoS 策略与跨平台支持等架构特性 [Macenski et al., 2022]。上述工作共同表明,边缘协同范式在降低控制延迟方面具有显著优势,但依赖特定基础设施VPN、DDS中间件)的部署方式在资源受限的嵌入式场景中仍存在适用性问题。
微服务化架构的引入。 微服务架构Microservice Architecture)将单体应用分解为一组职责单一、可独立部署的服务单元,通过轻量级接口进行协作,已在 Web 工程领域得到广泛验证。Xia 等人首次将微服务架构引入云机器人系统,利用 Docker 容器将机器人功能模块化,并通过 Kubernetes 进行编排管理,在 Visual SLAM 任务中验证了微服务化带来的灵活性与可扩展性提升 [Xia et al., 2018]。微服务化架构的引入为本文将控制功能Control Service)与流媒体功能Stream Service)解耦部署于 UAV 边缘端提供了直接的设计范式参照。
综合来看,机器人控制架构经历了"集中式→分布式→云化→边缘协同→微服务化"的完整演进链条。本文所构建的三层异构协同架构正是在吸收边缘协同与微服务化两大范式优势的基础上,针对动物机器人特殊约束进行适应性设计的产物。

\subsection{异构协同架构的任务分配理论}
任务卸载Task Offloading)是边缘协同架构的核心决策机制,指将部分或全部计算任务从资源受限的本地节点迁移至算力更强的远端节点边缘服务器或云端)执行的过程。本节梳理任务卸载的理论框架,并重点分析动态卸载策略在动物机器人控制场景中的适用性局限。
任务卸载的分类框架。 根据决策时机与灵活程度,任务卸载策略通常分为三类：

静态卸载：在系统部署阶段预先确定任务执行位置,运行时不再调整。该策略延迟确定性高,但缺乏对网络状态变化的适应能力；
动态卸载：在运行时根据网络质量、计算负载等实时指标动态调整任务执行位置。该策略自适应性强,但引入调度决策本身的延迟开销,且存在任务中途切换导致状态不一致的风险；
混合卸载：对不同类型任务采用差异化策略——实时性强的任务采用静态分配,计算密集型任务采用动态卸载。该策略在延迟确定性与资源利用率之间寻求平衡。

代表性工作分析。 Li 等人提出的 EARNEdge Accelerated Robot Navigation)框架在机器人本地保守规划器与边缘端高性能规划器之间动态切换,通过模型预测切换Model-Predictive Switching,MPS)最大化边缘加速收益 [Li et al., 2025]。Satpure 等人针对 UAV 场景,将模型预测控制MPC)控制器基于 Kubernetes 和 Docker 容器化部署于边缘服务器,控制频率从机载20 Hz提升至100 Hz [Satpure et al., 2022]。Suganya 等人提出面向 UAV 的动态任务卸载框架DTOE-AOF),以地面站作为系统枢纽,通过 Wi-Fi 和5G与边缘节点通信,边缘节点在地面站和 UAV 之间充当中介,实现了三层数据流的协同 [Suganya et al., 2024]。
动态卸载对动物机器人刺激控制的潜在危害。 上述动态卸载方案均以机械执行器为控制对象——舵机、推进器等设备对指令的时序偏差具有一定的容忍裕度。然而,动物机器人的神经电刺激控制对时序确定性要求显著更高：大鼠内侧前脑束的奖励刺激有效窗口通常在数十至数百毫秒量级,若因动态卸载决策引入不可预测的延迟抖动,将导致刺激-行为关联弱化乃至训练效果丧失。此外,生物体行为本身具有不确定性,动态卸载在切换执行节点时若出现中间状态不一致,可能引发指令丢失或重复刺激,对实验的可重复性产生严重干扰。
因此,本文在任务调度策略上采用静态-动态混合方案：神经刺激控制指令的下发任务固定静态)保留于 UAV 边缘端本地执行,以保证刺激时序的确定性；视觉处理、路径规划等计算密集型任务则依据链路质量指标进行动态卸载决策。这一策略的理论必要性正源于上述对纯动态卸载方案的局限性分析,其具体算法设计将在第4章详细阐述。
\subsection{面向嵌入式设备的微服务部署}
将微服务架构部署于嵌入式边缘设备,面临与云端部署显著不同的约束：计算资源有限Jetson Orin 等嵌入式平台 RAM 通常为8至32 GB,远低于云服务器)、存储空间受限、且须保证实时响应而非仅追求吞吐量。本节梳理该方向的代表性工作,并指出其对本文场景的适用性边界。
安全关键场景的微服务设计。 Schrick 等人将微服务架构应用于安全关键移动机器人控制,提出基于微服务和层次有限状态机HFSM)的控制框架,通过微服务化减少集中式逻辑,显著提升系统并发性、可中断性和容错能力 [Schrick et al., 2024]。其核心贡献在于将"单点失效"风险分散至各微服务单元,单个服务的故障不影响其他服务的正常运行,这一容错思想对本文的服务生命周期设计具有重要参考价值。
嵌入式系统的微服务框架。 针对嵌入式系统资源受限特点,IEEE COMPSAC 2021的研究提出了适用于嵌入式平台的轻量级微服务框架,设计了嵌入式服务模型、轻量通信协议和服务注册发现流程 [Placeholder\_COMPSAC2021],为在 Jetson Orin 等设备上部署微服务提供了可行性验证。Satpure 等人在 UAV 场景中验证了基于 Kubernetes 和 Docker 的边缘部署方案 [Satpure et al., 2022],证明容器化编排技术可在计算资源受限的边缘节点上稳定运行。
现有方案的适用性边界。 综合上述研究可以发现,现有嵌入式微服务方案均未针对超低带宽下行链路如大鼠背包端约200 kbps 的蓝牙通道)进行服务接口设计,其服务间通信均假设充足的局域网带宽环境。将上述方案直接应用于本文场景时,Control Service 与背包端之间的通信设计将成为性能瓶颈,标准微服务通信协议的协议头开销将严重挤占有限的控制信道带宽。这一未解决的问题,正是本文提出 mice-gRPC 轻量化通信协议的直接动机,相关协议设计将在第5章详细阐述。

\section{实时通信协议的选型依据}
本节从协议选型的工程视角出发,系统对比机器人控制领域主流通信协议的性能特征,重点分析 gRPC 协议的技术原理,并指出现有协议在资源受限环境下的优化缺口,为本文通信协议设计提供理论依据
\subsection{机器人控制通信协议演进}
机器人控制通信协议经历了从简单无线链路到标准化中间件、再到轻量化应用层协议的多代演进,不同协议在延迟、带宽、可靠性与部署复杂度等维度上呈现出显著差异。本节从以上四个维度对主流协议进行横向对比见表2.1),为本文的协议选型提供定量依据。
蓝牙/BLE 协议在动物机器人领域被广泛应用于背包端与地面站之间的无线通信,其优势在于功耗低、配对简单、硬件成本低廉。然而,Zhang 等人的脑-脑接口大鼠控制实验实测显示,蓝牙控制指令延迟达155至495 ms [Zhang et al., 2019],且单通路带宽约200 kbps,无法支持控制指令与视频流的并行传输；蓝牙有效通信距离通常不超过10 m,在 Rat-UAV 场景中 UAV 飞行半径大幅超出覆盖范围 [PlaceholderBLE2017]。
DDS/ROS 2 是机器人领域应用最广泛的中间件通信标准。ROS 2 采用 DDSData Distribution Service)实现去中心化点对点通信,支持丰富的 QoS 策略配置。然而,Maruyama 等人的实测表明,ROS 2 相比底层 DDS 可产生约50\%的延迟附加开销 [Maruyama et al., 2016],且在树莓派等嵌入式设备上无法保证1 ms 通信周期 [PlaceholderROS2Realtime2023],对本文小于11 ms 端到端延迟目标而言裕度不足。
MQTT 是面向物联网场景的轻量级发布-订阅协议,协议头开销极小最小2字节),适合带宽极受限的传感器网络。但 MQTT 默认基于 TCP,缺乏原生的请求-响应语义和流式传输能力,难以直接用于需要双向同步控制的机器人场景 [Profanter et al., 2019]。
gRPC 基于 HTTP/2 和 Protocol Buffers,在微服务通信场景中展现出显著的延迟优势。对比测试表明,gRPC 响应时间比 REST API 快约77\% [PlaceholderGrpcRest2024],其多路复用特性支持单连接上的并发 RPC 调用,Protocol Buffers 序列化消息体积约为 JSON 的1/2至1/3 [Popic et al., 2016],在延迟、带宽利用率和接口规范化方面均具有明显优势。

表2.1 机器人控制主流通信协议横向对比
协议名称典型延迟带宽占用可靠性机制适用场景主要局限蓝牙/BLE155–495 ms [Zhang et al., 2019]约200 kbps链路层重传短距低功耗IoT延迟高、距离短、单通路DDS/ROS 2基础延迟+50\%开销 [Maruyama et al., 2016]较高协议头重)QoS多级保障工业机器人、多机协同嵌入式设备开销大MQTT<10 ms局域网)极低2字节头)QoS 0/1/2级物联网传感器网络缺乏双向同步语义gRPC<5 ms局域网)[PlaceholderGrpcRest2024]低Protobuf序列化)HTTP/2流控微服务RPC调用远距离无帧保护机制REST/HTTP比gRPC慢约77\% [PlaceholderGrpcRest2024]高JSON/XML)应用层重试Web API延迟高、序列化开销大

综合上述对比可以看出,gRPC 在延迟性能和序列化效率方面具有最显著的综合优势,是本文地面站-边缘端间通信的基础选型。然而,gRPC 标准实现在远距离低带宽场景下的可靠性保障机制仍存在明显不足,这将在2.3.3节进一步分析。

\subsection{gRPC协议性能分析与适用性}
gRPCGoogle Remote Procedure Call)是 Google 开源的高性能 RPC 框架,基于 HTTP/2 传输层和 Protocol BuffersProtobuf)序列化协议构建,已成为微服务通信领域的事实标准之一。本节从协议原理和性能特征两个层面进行深入分析。
HTTP/2 多路复用原理。 gRPC 基于 HTTP/2 协议,其核心优势在于多路复用Multiplexing)机制：在单条 TCP 连接上,HTTP/2 通过流Stream)机制支持多个并发请求的交错传输,彻底消除了 HTTP/1.1 中因请求串行化导致的队头阻塞Head-of-Line Blocking)问题。对于本文场景,这意味着控制指令 RPC 调用与状态上报 RPC 调用可以在同一 TCP 连接上并发执行,无需建立多条独立连接,显著降低了连接建立的开销。
Protocol Buffers 序列化机制。 gRPC 默认使用 Protocol BuffersProtobuf)作为接口定义语言IDL)和序列化格式。Protobuf 采用变长编码Varint)对整型字段进行压缩,序列化后的消息体积约为等效 JSON 的1/3至1/2,序列化速度显著快于 JSON 和 XML [Popic et al., 2016]。在带宽约为200 kbps 的背包端蓝牙链路上,消息体积的压缩对于提升有效控制帧传输速率具有直接意义。
Unary RPC 与 Server Streaming RPC 的选择逻辑。 gRPC 提供四种通信模式：Unary RPC单次请求-响应)、Server Streaming RPC服务端流)、Client Streaming RPC客户端流)和 Bidirectional Streaming RPC双向流)。在本文架构中,控制指令下发对实时性要求高、消息体积小,适合采用 Unary RPC 模式——单次 RPC 调用即可完成指令发送与执行确认的闭环,延迟开销最小；视频流数据具有持续性和高带宽特征,适合采用 Server Streaming RPC 模式——服务端持续推送帧数据,客户端无需反复发起请求,减少了连接建立开销。
嵌入式设备上的性能评估。 在 Jetson Orin 等 ARM64 嵌入式平台上,gRPC Python 实现的单次 Unary RPC 往返延迟在局域网环境下通常低于5 ms [PlaceholderGrpcRest2024],满足本文小于11 ms 端到端控制延迟的设计目标。相比之下,ROS 2 在同类嵌入式设备上的延迟附加开销约为底层 DDS 的50\% [Maruyama et al., 2016],难以满足相同的实时性要求。这一对比从协议性能角度支撑了本文选择 gRPC 而非 ROS 2 作为边缘端服务通信框架的技术决策。
\subsection{资源受限环境下的协议优化方向}
尽管 gRPC 在通用微服务场景中具有显著优势,但其标准实现在面向大鼠背包端超低带宽、远距离无线链路等资源受限环境时,仍存在若干尚未解决的工程问题。本节系统梳理相关优化方向,并指出现有研究的局限。
帧压缩与协议头精简。 HTTP/2 协议帧头部固定开销为9字节,Protobuf 消息在经过 gRPC 封装后,控制帧总大小通常在60至80字节范围内。在200 kbps 的蓝牙链路上,若以20 Hz 频率下发控制指令,仅控制帧即占用约12.8 kbps80字节x20 Hzx8 bit),考虑视频流、状态心跳等并发传输需求,带宽竞争将显著影响控制帧的及时送达。因此,针对控制帧的极致轻量化设计——压缩 Protobuf 字段、减少不必要字段填充——是资源受限场景下的重要优化方向。
优先级队列与控制帧保护。 标准 gRPC 基于 HTTP/2 的流量控制机制对所有 RPC 调用一视同仁,不区分控制帧与数据帧的传输优先级。当视频流Stream Service)与控制指令Control Service)共享同一网络链路时,大体量的视频帧极易占满带宽窗口,导致控制帧排队等待,引入不可预测的附加延迟。引入双优先级队列机制——控制帧进入高优先级队列、视频帧进入低优先级队列——是解决该问题的关键方向,但目前尚无面向动物机器人通信场景的实现方案。
ARQ 重传机制与坏帧保护。 标准 gRPC 依赖 TCP 的可靠传输保障,TCP 的重传机制虽然保证了最终可靠性,但在远距离无线场景中信号衰减导致帧错误率上升),TCP 的拥塞控制算法会将重传触发的超时误判为网络拥塞,进而降低发送窗口,造成吞吐量的大幅下降。面向无线链路的 ARQAutomatic Repeat reQuest)机制在应用层提供帧级错误检测与重传,可在不触发 TCP 拥塞控制的前提下快速恢复丢失帧,对于保障远距离控制指令的可靠性至关重要。
单通路多路复用。 大鼠背包端受体积与功耗限制,通常仅能搭载单一蓝牙通信模块,形成单通路约束。如何在单条蓝牙链路上同时传输控制指令、视频流和状态心跳三类数据,并保证各类数据的实时性需求,是现有 gRPC 框架未涉及的工程问题。时分多路复用Time-Division Multiplexing,TDM)结合帧调度算法是解决该问题的主要技术路径。
综合上述分析,现有 gRPC 标准协议在远距离低带宽场景中存在协议头开销偏重、无控制帧优先保护机制、缺乏应用层坏帧保护以及不支持单通路多优先级复用四个核心局限。这些局限在通用 Web 微服务场景中影响有限,但在大鼠背包端200 kbps 蓝牙链路和15 m以上导航距离的约束下将被显著放大,成为制约系统实时性与可靠性的关键瓶颈。这一理论缺口正是本文在第5章提出 mice-gRPC 协议的核心动机。


\section{动物机器人信息传输系统现状与问题}
本节以问题驱动的视角系统审视现有动物机器人信息传输系统的架构模式与量化性能,通过与 Rat-UAV 协同导航任务需求的对比,揭示现有系统在延迟、距离、带宽三个维度上的结构性局限,并从工程可行性角度论证引入分布式协同架构的必要性。
\subsection{现有动物机器人通信系统的共性架构}
尽管动物机器人控制算法经历了从简单遥控到闭环自主导航的显著进步,但其底层信息传输系统的架构模式自2002年以来几乎未发生根本性变化,普遍沿用"地面PC集中处理 + 蓝牙点对点传输"的单体通信模式。
典型系统的架构分析。 Feng 等人于2007年设计了基于蓝牙无线微刺激器的大鼠遥控系统,采用 C8051 微处理器驱动的蓝牙模块实现地面端与背包端的点对点通信,地面端 PC 承担全部信号处理与指令生成职责 [Feng et al., 2007]。Wang 等人于2015年提出的视觉线索引导框架将摄像头采集图像经无线传输至 PC 端进行目标检测,再将决策结果转化为刺激序列下发至背包端 [Wang et al., 2015]——其数据流本质仍为"背包传感→无线→PC处理→无线→背包执行"的单环路集中式模式。
Zhang 等人于2019年实现的脑-脑接口系统同样采用蓝牙作为唯一通信链路,实测控制指令延迟达155至495 ms [Zhang et al., 2019],尽管其应用场景人脑意念控制大鼠)与本文不同,但所采用的通信架构与性能瓶颈完全一致。Zheng 等人于2023年提出的基于有限状态机的大鼠自主导航系统在算法层面实现了重要突破,但其控制架构仍为单板嵌入式系统加集中式处理 [Zheng et al., 2023]。即便是最新的 Rat-UAV Navigation 工作,其软件控制架构依然以地面端 PC 为中心,UAV 仅充当视觉传感器的搭载平台,尚未承担边缘计算节点的功能 [Zheng et al., 2024]。
上述分析清晰表明,在23年的发展历程中,动物机器人控制系统在算法创新层面取得了丰硕成果,但在系统架构层面从未突破集中式单体模式的范式约束。

\subsection{现有系统的通信局限与研究空白}
以 Rat-UAV 协同导航任务为参照基准,本节从延迟、距离、带宽三个维度对现有蓝牙集中式方案进行量化对比分析,揭示其在当前任务场景下的结构性局限见表2.2)。
延迟维度。 现有系统的蓝牙通信延迟为155至495 ms [Zhang et al., 2019]。神经电刺激控制的有效时序要求通常在数十毫秒量级——奖励刺激需在大鼠完成目标行为后及时施加以建立正强化关联；若刺激延迟超过行为发生后约500 ms,条件化效果将显著下降 [PlaceholderNeurostim]。当前蓝牙方案的延迟上限495 ms)已接近神经刺激有效窗口的极限,且存在较大的延迟抖动,在 Rat-UAV 场景中 UAV 飞行引入的额外信号衰减将进一步恶化延迟表现。
距离维度。 蓝牙的有效通信距离在障碍物干扰环境中通常为3至10 m [PlaceholderBLE2017],而 Rat-UAV 协同导航场景中 UAV 的典型飞行半径在数十至百余米量级。两者之间存在一个数量级以上的距离差距,使现有蓝牙方案在 Rat-UAV 场景中从根本上不可用。
带宽维度。 蓝牙4.x 的物理层净荷吞吐量约为200 kbps,Rat-UAV 场景中需同时传输的数据包括：控制指令流约10–20 kbps)、摄像头视频流即使压缩后也需要数百 kbps)以及状态心跳约5 kbps)。三类数据的带宽需求之和远超蓝牙单通路的承载能力,导致现有系统不得不在控制指令与视频流之间作出取舍,无法实现两者的并行传输。

表2.2 现有蓝牙方案与 Rat-UAV 任务需求对比
指标现有蓝牙方案Rat-UAV 任务需求差距控制延迟155–495 ms [Zhang et al., 2019]<50 ms神经刺激时序要求)差距3–10倍通信距离3–10 m [PlaceholderBLE2017]>50 mUAV飞行半径)差距5–15倍可用带宽~200 kbps单通路)>500 kbps控制+视频+状态)差距>2.5倍多路复用不支持控制/视频/状态并行传输结构性缺失

\subsection{引入分布式协同架构的必要性论证}
面对上述量化局限,一个自然的问题是：能否直接借鉴工业机器人或云机器人领域已有的分布式架构方案加以解决？本节对两类候选方案的适用性进行逐一论证,得出必须针对动物机器人特殊约束重新设计的结论。
工业机器人分布式方案不可直接照搬。 工业机器人分布式控制方案如 FogROS2、CIRPO 等)均以确定性机械执行器为控制对象：电机、液压机构等设备对指令延迟的容忍度在数十至数百毫秒量级,且其执行结果具有精确的物理确定性。动物机器人则面临三项根本性差异：其一,生物体不确定性——大鼠个体在神经接口效果、疲劳程度、情绪状态等方面存在显著差异,控制系统须具备实时闭环反馈能力,不可采用开环批量指令模式；其二,背包端资源极度受限——大鼠背包受体积和重量限制通常不超过10 g [PlaceholderBackpackWeight]),仅能搭载超低功耗嵌入式单元,无法运行完整的 DDS 中间件或 Kubernetes 代理；其三,神经刺激对延迟确定性要求更高——工业执行器对50 ms以内的延迟变动通常不敏感,但神经刺激若在有效时序窗口外施加,将直接导致强化学习训练失效,系统须保证刺激指令的确定性时延上界。上述三点差异共同决定了工业机器人分布式方案无法直接迁移至动物机器人场景。
云机器人方案FogROS2)不可直接采用。 FogROS2 等云机器人平台通过公有云 VPN 实现计算节点的透明迁移,其网络往返延迟通常超过50 ms [Ichnowski et al., 2023],远不满足本文小于11 ms 控制延迟的设计目标。此外,FogROS2 依赖公有云基础设施,在 UAV 飞行于信号覆盖稀疏区域时将面临完全断连风险,且缺乏离线降级机制——一旦云端连接中断,控制系统将无法切换至本地模式保证基本导航功能。对于须在野外复杂环境中执行长时任务的 Rat-UAV 系统而言,这一脆弱性不可接受。
结论：必须针对动物机器人特殊约束设计专用分布式协同架构。 综合上述分析,动物机器人自动导航场景面临的约束组合——背包端超低带宽约200 kbps)、生物体不确定性要求的实时闭环、UAV 飞行的弱网/远距离通信、神经刺激严苛的时序确定性需求——共同构成了一个现有工业和云机器人分布式方案均无法覆盖的特殊约束空间。本文因此提出针对上述约束量身定制的三层异构协同控制架构,并在此基础上设计面向弱网的低时延任务调度算法WLTS,第4章)和面向远距离导航的 mice-gRPC 通信协议第5章),以系统性地突破上述瓶颈。

\section{本章小结}
本章围绕面向动物机器人地面站-边缘端分布式协同控制系统所需的三类核心技术基础,系统梳理了分布式控制架构设计方法、实时通信协议选型依据以及现有动物机器人信息传输系统的局限与研究空白。表2.3对三个方向的核心结论作综合性汇总。

表2.3 本章文献综述核心结论汇总
研究方向现有代表性工作核心局限本文对应解决方案分布式架构FogROS2 [Ichnowski et al., 2023]、微服务云机器人 [Xia et al., 2018]面向工业/云场景,未适配背包端超低带宽与离线降级需求三层异构协同控制架构第3章)任务调度EARN [Li et al., 2025]、DTOE-AOF [Suganya et al., 2024]纯动态卸载引入刺激时序不确定性,不适用于神经刺激控制弱网低时延静态-动态混合调度算法 WLTS第4章)通信协议蓝牙 [Zhang et al., 2019]、标准 gRPC [PlaceholderGrpcRest2024]蓝牙延迟高/距离短；gRPC 无控制帧优先保护与坏帧恢复机制mice-gRPC 轻量化扩展通信协议第5章)

上述三类问题相互关联、层层递进：架构设计确定了任务调度的边界条件,任务调度的正确执行依赖通信协议提供可靠的低延迟数据链路,而通信协议的设计又须在架构约束背包端资源限制、单通路通信)下进行优化。第3章将在本章理论基础之上,详细阐述面向动物机器人的三层异构协同控制架构的系统设计。

\chapter{第3章 面向弱网环境的低时延任务调度算法}

\section{引言}

\section{问题建模与延迟分析}

\subsection{端到端控制延迟的构成分析}


\subsection{弱网场景下的延迟波动建模}


\subsection{任务调度问题形式化定义}


\section{XXX任务动态调度算法设计}
\subsection{链路质量感知模块}


\subsection{静态-动态混合任务调度策略}

\subsection{容错降级机制}

\section{算法复杂度与理论分析}


\section{实验验证}
\subsection{实验设置}

\subsection{延迟性能实验}

\subsection{与现有方案对比}



\section{本章小结}



\chapter{第4章 面向远距离导航的 mice-gRPC 通信协议}

\section{引言}



\section{问题分析与协议需求}
\subsection{动物机器人通信链路特性分析}

\subsection{标准gRPC在此场景的局限}

\subsection{mice-gRPC的设计目标与约束}


\section{mice-gRPC 协议设计}

\subsection{轻量化控制帧格式}

\subsection{坏帧校验机制}

\subsection{优先级重传机制}


\section{协议实现}
\subsection{gRPC拦截器实现}

\subsection{帧编解码模块}

\subsection{优先级队列调度器}

\section{实验验证}
\subsection{实验设置}

\subsection{延迟与帧成功率测试}

\subsection{带宽占用对比}

\section{本章小结}


\chapter{第5章 面向动物机器人的三层异构协同控制架构设计}
\section{系统架构需求分析}
\subsection{Rat-UAV协同范式分析}
\subsection{三层架构的功能需求}
\subsection{异构平台约束分析}

\section{三层协同架构总体设计}
\subsection{架构设计原则}
\subsection{三层协同架构模型}
\subsection{各层职责划分与交互协议}

\section{微服务化部署设计}
\subsection{服务划分原则与边界定义}
\subsection{mice-gRPC服务接口设计}
\subsection{服务生命周期与容错机制}

\section{容器化部署方案}
\subsection{Docker多架构镜像构建}
\subsection{Compose编排与健康检查}

\section{本章小结}


\chapter{第6章 系统集成与实验结果分析}

\section{任务场景描述}
\subsection{动物机器人自动导航任务定义}


\subsection{关键技术指标体系}


\section{系统部署与运行效果展示}


\subsection{三层架构部署配置}



\subsection{系统运行效果}


\subsection{平台在任务中的关键作用分析}



\section{综合性能对比实验}
\subsection{与现有动物机器人系统对比}

```

\subsection{架构模式对比}


\subsection{通信协议对比}

\section{系统局限性讨论}

\section{本章小结}




\chapter{第7章 总结与展望}
\section{全文总结}
动物机器人作为一种新型的生物-机械融合系统,在灾害搜救、环境监测等领域具有独特的应用价值。然而,现有控制系统普遍面临控制延迟过高、集中式架构导致单点故障、缺乏分布式部署能力等问题,制约了动物机器人技术的进一步发展。针对这些问题,本文提出了一种面向动物机器人自动导航的地面站-边缘端协同控制架构,并完成了系统的设计、实现和验证工作。

本文的主要研究工作和贡献总结如下：

1)提出了面向动物机器人的地面站-边缘端协同控制架构

本文设计了一种三层分布式架构,将系统划分为地面站层、边缘端层和硬件层。边缘端部署Control Service和Stream Service两个微服务,分别负责控制指令处理和视频流传输；地面站部署导航算法和用户界面,通过gRPC协议与边缘端通信。这一架构实现了计算任务的合理分配：延迟敏感型任务刺激控制、传感采集)在边缘端就近处理,计算密集型任务视觉处理、路径规划)在地面站执行。相比传统的集中式架构,分布式设计提高了系统的可靠性和扩展性,支持地面站与边缘端的物理分离部署,为多动物并发控制奠定了基础。

2)实现了基于gRPC的低延迟控制协议

本文针对动物机器人控制场景的实时性需求,设计并实现了基于gRPC Unary RPC的控制协议。通过Protocol Buffers的高效序列化消息体积比JSON小3-10倍)、HTTP/2的持久连接和多路复用、以及串口通信的优化策略连接池管理、指令格式优化),系统实现了8-11ms的端到端控制延迟。与现有动物机器人系统BBI系统155ms、Rat-UAV系统500ms)相比,延迟降低了90\%以上；与REST API相比,延迟降低了约70\%。这一性能提升对于需要精确时序控制的神经科学实验具有重要价值。

3)实现了基于gRPC Server Streaming的视频流传输优化

本文采用gRPC Server Streaming机制实现视频流从边缘端到地面站的传输。系统支持JPEG和H.264两种编码方案,可根据网络条件和硬件能力自适应切换。通过帧缓冲管理双缓冲机制保证读写分离)和流控机制背压控制避免网络拥塞),在WiFi网络环境下实现了10fps稳定帧率、50-90ms端到端延迟、13.9Mbps带宽占用、零丢帧的传输性能。针对Windows环境下摄像头启动慢的问题,通过指定DirectShow后端将启动时间从超过10秒优化至2.9秒,性能提升71\%。

4)设计了基于适配器模式的渐进式系统改造方法

本文设计了RemoteBagVJ和RemoteCameraVJ两个适配器类,封装gRPC通信细节,对外提供与原有单体架构完全一致的接口。这一设计遵循"开闭原则",使得上层业务代码无需任何修改即可切换到分布式架构。实验验证表明,适配器与原接口的兼容性达到100\%,包括背包控制、摄像头管理、视频录制、传感器读取等全部功能。透明切换机制支持通过配置文件在本地模式和远程模式之间灵活切换,降低了系统改造的风险和成本。

5)实现了面向异构平台的统一部署方案

本文基于Docker容器技术,设计了支持x86\_64和ARM64双架构的统一部署方案。通过Docker Buildx多架构镜像构建,同一个镜像标签可以在Windows PC和NVIDIA Jetson等不同平台上自动选择对应版本运行。Docker Compose服务编排实现了多服务的一键部署,健康检查机制保证了服务的可用性监控和自动恢复。这一方案简化了在异构边缘设备上的部署流程,提高了系统的可移植性。

6)完成了系统的全面实验验证

本文在单机环境场景A)和双机网络环境场景B)下对系统进行了全面的测试和验证。功能测试验证了Control Service和Stream Service的独立功能以及跨服务的协同工作能力,离线视频导航集成测试验证了微服务架构对单体架构核心功能的完全兼容。性能测试表明系统达到了所有设计指标：控制延迟8-11ms目标<50ms)、视频帧率10fps目标≥10fps)、传输带宽13.9Mbps目标<50Mbps)。对比实验证明了gRPC相比REST在延迟和吞吐量方面的显著优势,以及本文系统相比现有动物机器人系统在延迟方面的数量级改进。
\section{工作展望}
基于本文工作的基础和存在的局限性,未来可以在以下方向开展深入研究：

1)真实环境部署与验证

下一步工作的首要任务是将系统部署到真实的动物机器人实验环境中进行验证。这包括：在NVIDIA Jetson Orin等边缘设备上测试系统的性能和功耗；连接真实的背包硬件验证串口通信的可靠性；在实际的大鼠导航实验中评估系统对实验效果的影响。真实环境的验证将为系统的实用化奠定基础。

2)H.264/NVENC硬件编码集成

充分利用Jetson平台的NVENC硬件编码能力,实现高效的H.264视频压缩。这将显著降低视频传输的带宽需求,使系统能够适应更广泛的网络环境。同时,可以研究自适应码率控制算法,根据网络状况动态调整视频质量,实现带宽与质量的最优权衡。

3)多动物协同控制扩展

扩展系统架构以支持多动物并发控制。这需要解决以下技术问题：设计基于服务发现机制如Consul、etcd)的动态服务注册；实现负载均衡策略以合理分配计算资源；设计多动物状态同步协议以支持协同导航实验。多动物控制能力将大幅提升系统的实验效率和应用价值。

4)边缘智能集成

探索在边缘端部署轻量级导航算法的可能性。随着边缘设备算力的提升,部分视觉处理和决策功能可以下沉至边缘端执行,进一步降低系统延迟。可以研究模型压缩、知识蒸馏等技术,将复杂的深度学习模型适配到资源受限的边缘设备上。边缘智能的集成将实现更低延迟、更高自主性的动物机器人控制。

5)云边端协同架构

将当前的地面站-边缘端两层架构扩展为云-边-端三层架构。云端负责大规模数据存储、模型训练和跨实验分析；边缘端负责实时控制和数据采集；端侧动物背包)负责刺激执行和传感感知。这种架构将支持更大规模的动物机器人实验,并为数据驱动的神经科学研究提供基础设施支撑。



% 本模板根据浙江大学研究生院编写的《浙江大学研究生学位论文编写规则》~\cite{zjugradthesisrules},
% 在原有的 zjuthesis 模板~\cite{zjuthesis}基础上开发而来。

% 本模板的本科生版本\cite{zjuthesisrules}得到了浙江大学本科生院老师的支持与审核,
% 已经在本科生院网上公示。
% 但当前的研究生版本并未经过研究生院老师的审核,
% 同学们使用时要注意对照模板与要求,
% 切不可盲目使用。

% 作者本人并未编写过浙江大学研究生毕业论文,
% 所以不清楚具体要求。
% 如果有热心同学愿意帮忙,
% 可以替我联系相关老师,我会配合审核并修改代码。

% 如果你在Overleaf上编译本模板,请注意如下事项：

% \begin{itemize}
%    \item 删除根目录的 ``.latexmkrc'' 文件,否则编译失败且不报任何错误
%    \item 字体有版权所以本模板不能附带字体,请务必手动上传字体文件,并在各个专业模板下手动指定字体。
%        具体方法参照 GitHub 主页的说明。
 %   \item 当前的Overleaf默认使用TexLive 2017进行编译,但一些伪粗体复制乱码的问题需要TexLive 2019版本来解决。
  %      所以各位同学可以在Overleaf上编写论文时务必切换到TexLive 2019或更新版本来编译,以免产生查重相关问题。
   %     具体说明参照 GitHub 主页。
%\end{itemize}


%\section{节标题}

% 我们可以用includegraphics来插入现有的jpg等格式的图片,
% 如\autoref{fig:zju-logo}所示。

% \begin{figure}[htbp]
%     \centering
%     \includegraphics[width=.3\linewidth]{figure/img/rviz.png}
%     \caption{\label{fig:zju-logo}Rviz}
% \end{figure}


% \subsection{小节标题}


% \par 如\autoref{tab:sample}所示,这是一张自动调节列宽的表格。

% \begin{table}[htbp]
%     \caption{\label{tab:sample}自动调节列宽的表格}
%     \begin{tabularx}{\linewidth}{c|X<{\centering}}
%         \hline
%         第一列 & 第二列 \\ \hline
%         xxx & xxx \\ \hline
%         xxx & xxx \\ \hline
%         xxx & xxx \\ \hline
%     \end{tabularx}
% \end{table}


% \par 如\autoref{equ:sample},这是一个公式

% \begin{equation}
%     \label{equ:sample}
%     A=\overbrace{(a+b+c)+\underbrace{i(d+e+f)}_{\text{虚数}}}^{\text{复数}}
% \end{equation}

% \chapter{另一章}


% \begin{figure}[htbp]
%     \centering
%     \includegraphics[width=.3\linewidth]{example-image-a}
%     \caption{\label{fig:fig-placeholder}图片占位符}
% \end{figure}

% \chapter{再一章}

% \par 如\autoref{alg:sample},这是一个算法

% \begin{algorithm}[H]
%     \begin{algorithmic} % enter the algorithmic environment
%         \REQUIRE $n \geq 0 \vee x \neq 0$
%         \ENSURE $y = x^n$
%         \STATE $y \Leftarrow 1$
%         \IF{$n < 0$}
%             \STATE $X \Leftarrow 1 / x$
%             \STATE $N \Leftarrow -n$
%         \ELSE
%             \STATE $X \Leftarrow x$
%             \STATE $N \Leftarrow n$
%         \ENDIF
%         \WHILE{$N \neq 0$}
%             \IF{$N$ is even}
%                 \STATE $X \Leftarrow X \times X$
%                 \STATE $N \Leftarrow N / 2$
%             \ELSE[$N$ is odd]
%                 \STATE $y \Leftarrow y \times X$
%                 \STATE $N \Leftarrow N - 1$
%             \ENDIF
%         \ENDWHILE
%     \end{algorithmic}
%     \caption{\label{alg:sample}算法样例}
% \end{algorithm}